\section{Structural LEDH and Surrogate HMC}
\label{sec:bf-structural-ledh-ab-campaign}

The experiment completed on 24 October 2026 asks whether an LEDH proposal
defined on the declared innovations can support parameter estimation when
part of the state transition is deterministic.  It tests a linear model with
an exact likelihood and a nonlinear model with a deterministic completion.
Both four-chain HMC runs met the prescribed retained-sample $\widehat R<1.01$
criterion.  The linear model also passed a comparison of posterior means with
a refined exact-likelihood calculation.  The nonlinear run samples a fixed
finite-particle posterior; agreement with its exact model posterior remains
unestablished.  These are different conclusions, and the construction below
explains why.

Here \emph{campaign Model A} and \emph{campaign Model B} name the two software
test fixtures.  These letters do not identify Worked Examples A and B earlier
in this chapter.  The experiment extends the five-model case study in
Chapter~\ref{ch:bf-surrogate-hmc-five-models}; its data and configurations are
separate from those earlier runs.

\subsection{The two models and their parameter distributions}

Write the state as $x_t=(m_t,k_t)^T$.  Each model has two state coordinates,
one observation, and one new stochastic innovation per date.  All innovations
$e_t$ and observation errors $\eta_t$ below are independent standard normal
variables, also independent of the initial state.
Campaign Model A is the companion representation of an AR(2) process:
\begin{align}
 m_t &= 0.35m_{t-1}-0.10k_{t-1}+\sigma e_t,
 & k_t &= m_{t-1}, \label{eq:bf-campaign-a-transition}\\
 y_t &= m_t+\tau\eta_t,
 & x_0 &\sim N(0,I_2).
\end{align}
The unknown parameters are $(\sigma,\tau)$, with generating values
$(0.25,0.15)$.  In coordinates $q=(\log\sigma,\log\tau)$ the prior is
\[
 q\sim N\!\left(
   \begin{bmatrix}\log 0.4\\\log 0.3\end{bmatrix},
   \operatorname{diag}(0.7^2,0.7^2)\right).
\]
There is no innovation in the second transition equation.  Nevertheless, its
predictive variance is generally positive because the preceding $m$ is
uncertain.  A Kalman filter handles the rank-one process covariance directly;
it provides the exact likelihood used at Model A's Metropolis endpoints.

Campaign Model B replaces the lag identity with a nonlinear completion:
\begin{align}
 m_t &= \rho m_{t-1}+\sigma e_t,
 & k_t &=0.55k_{t-1}+\beta\tanh(m_t),
 \label{eq:bf-campaign-b-transition}\\
 y_t &=m_t+k_t+0.30\eta_t,
 & x_0 &\sim N\!\left(0,\operatorname{diag}(0.25,0.20)\right).
\end{align}
The unknown parameters are $(\rho,\sigma,\beta)$, with generating values
$(0.7,0.25,0.8)$.  Define
$q=(\operatorname{atanh}\rho,\log\sigma,\log\beta)$ and use
\[
 q\sim N\!\left(
 \begin{bmatrix}0.5\\\log 0.4\\\log 0.7\end{bmatrix},
 \operatorname{diag}(0.8^2,0.7^2,0.6^2)\right).
\]
These are declared regularizing priors for synthetic experiments, not priors
attributed to an application paper.  Their densities are defined with respect
to $dq$.  Thus the HMC target in $q$ uses these normal densities directly;
one must not add a second Jacobian as though the same normal prior had been
defined on the physical parameters.  Credible intervals are computed after
transforming each retained draw back to the physical coordinates.

\subsection{Proposing an innovation and completing the state}

After the preceding filter reset, let $x_{t-1}^{(j)}$, $j=1,\ldots,N$, be
the equally weighted parent cloud.  Denote structural completion by
$F_\theta(x,e)$, using either
\eqref{eq:bf-campaign-a-transition} or
\eqref{eq:bf-campaign-b-transition}.  Conditional on a chosen parent, the
proposal is a density in the scalar $e$, not a bivariate Gaussian density on
$x_t$.  It has three components:
\begin{equation}
 r_j(e)=0.2\,\phi(e)+0.4\,r_{j,\mathrm{SP}}(e)
                       +0.4\,r_{j,\mathrm{LEDH}}(e),
 \label{eq:bf-structural-campaign-mixture}
\end{equation}
where $\phi$ is the standard normal density, $r_{j,\mathrm{SP}}$ is the
shared sigma-point Gaussian proposal, and $r_{j,\mathrm{LEDH}}$ is the shared
local-flow Gaussian proposal.  Both guided components observe the completed
state through $h(F_\theta(x_{t-1}^{(j)},e))$.  The required derivatives with
respect to $e$ are $\sigma$ in A and
\[
 \frac{\partial h(F_\theta(x,e))}{\partial e}
   =\sigma\{1+\beta[1-\tanh^2(m)]\}
\]
in B.  This derivative includes the effect of the innovation on the
deterministically completed coordinate.

The algorithm chooses a parent uniformly, chooses a mixture component, draws
its innovation, and computes $x_t=F_\theta(x_{t-1}^{(j)},e)$.
All three densities in \eqref{eq:bf-structural-campaign-mixture} are evaluated
at that same innovation and that same parent.  The incremental weight is
\begin{equation}
 w_t^{(i)}=
 \frac{g_\theta(y_t\mid F_\theta(x_{t-1}^{(j_i)},e_i))\phi(e_i)}
      {r_{j_i}(e_i)},
 \qquad
 \widehat\ell_F(q)=\sum_{t=1}^T
       \log\!\left\{\frac1N\sum_{i=1}^N w_t^{(i)}\right\}.
 \label{eq:bf-structural-campaign-weights}
\end{equation}
The uniform parent probability cancels between target and proposal on the
augmented space $(j,e)$.  Both numerator and denominator are densities with
respect to the same innovation measure $de$, so no ambient-state Jacobian is
needed.  Summing proposal densities over other parents would instead mix
different completion curves and would not compute this importance ratio.

Every propagated child satisfies its deterministic equation by construction.
The subsequent weighted-cloud reset approximates the unconditional filtering
distribution in state space; a reset point need not lie on a particular
previous parent's completion curve.  At the next propagation the conditional
structural identity is imposed again.  This distinction matters: structural
propagation avoids artificial process noise, but it does not make the
finite-cloud reset exact Bayesian filtering.

The structural adapter connects each model to the shared map builder.
It is implemented in \path{structural_completion_proposal_tf.py}.
The consumer in the shared
batched score recurrence selects this adapter for structural models, then
uses the shared weight, transport, moment-reset, and correction operations.
The state and all parameter tangents travel together through the recursion.
In particular,
\[
 D_qx_t=D_xF_\theta(x_{t-1},e)D_qx_{t-1}
          +D_eF_\theta(x_{t-1},e)D_qe
          +\partial_qF_\theta(x_{t-1},e),
\]
and the same chain rule differentiates proposal densities and weight
normalizations.  The production score is this analytical tangent calculation,
not a separate finite-difference filter run for each parameter.  Fixed draw
designs, component labels, and parent labels define the finite computation;
numerical validity and local branch behavior are checked separately.

The campaign also replaces the dense Cayley-matrix calculation in the shared
moment correction with an algebraically equivalent compact action.  If the
cosine basis is $U\in\mathbb R^{N\times2r}$, its skew core is $A$, and
$G=U^TU$, write the scaled skew generator as $K=UDU^T$.  The normalization
used by the original algorithm is retained:
\[
 D=\frac{\kappa A}{\sqrt{1+\operatorname{tr}(A^T G A G)}}.
\]
For the Cayley transform $O=(I-K/2)^{-1}(I+K/2)$, the correction needs
$O^TZ$, not the full $N\times N$ matrix.  Since
$O^T=I-(I+K/2)^{-1}K$ and
$(I+UDU^T/2)^{-1}U=U(I+DG/2)^{-1}$, it follows that
\[
 O^TZ=Z-U(I+DG/2)^{-1}DU^TZ.
\]
Only a $2r$-dimensional solve is required.  Retaining $G$, rather than
replacing it by the identity, preserves the actual finite cosine basis even
when its columns are not exactly orthogonal.  Tests compare the compact and
dense actions, parameter tangents, and mixed derivatives in FP32 and FP64,
including deliberately nonorthogonal columns, and compare full-filter values
and scores on both structural models.  This removes the dense intermediate;
the campaign does not provide a paired timing estimate of its speedup.

\subsection{The force, the endpoint, and NeuTra whitening}

Let $p_q$ be the stated prior density and define the finite-filter force
\[
 G_F(q)=\nabla_q[-\log p_q(q)-\widehat\ell_F(q)].
\]
The force supplies the momentum kicks of leapfrog.  The potential evaluated
in the Metropolis decision is
\begin{equation}
 U_E(q)=
 \begin{cases}
 -\log p_q(q)-\ell_{\mathrm{Kalman}}(q),&\text{Model A},\\
 -\log p_q(q)-\widehat\ell_F(q),&\text{Model B}.
 \end{cases}
 \label{eq:bf-structural-campaign-endpoints}
\end{equation}
For A the force is a surrogate for the gradient of an exact endpoint.  For B
the endpoint itself is the fixed finite-LEDH likelihood approximation.  Its
random design is frozen; this is not a pseudo-marginal sampler with a fresh
unbiased likelihood estimate at each proposal.  Metropolis correction cannot
remove B's finite-filter approximation error.

The proposal argument does not require the force to equal the endpoint
gradient.  A kick $(q,p)\mapsto(q,p-\epsilon G_F(q))$ and a drift
$(q,p)\mapsto(q+\epsilon M^{-1}p,p)$ are triangular shears with determinant
one.  Symmetric composition and momentum reversal give the usual reversible,
volume-preserving proposal in exact arithmetic.  Acceptance using
$\min\{1,\exp[-H_E(q',p')+H_E(q,p)]\}$ therefore corrects this proposal
against the declared endpoint.  Numerical checks remain necessary for the
implemented finite-precision calculation; the argument does not identify
the two endpoints in \eqref{eq:bf-structural-campaign-endpoints}.

Model A uses affine tuning.  Model B uses a frozen NeuTra map $q=\mathcal T(z)$
trained against its own finite-filter posterior.  Its transformed endpoint
and force include the map's change of volume:
\begin{align}
 U_Z(z)&=U_E(\mathcal T(z))-\log|\det D\mathcal T(z)|,\\
 G_Z(z)&=D\mathcal T(z)^T G_F(\mathcal T(z))
             -\nabla_z\log|\det D\mathcal T(z)|.
\end{align}
The existing author-style IAF training implementation and the shared frozen-map
HMC bindings are reused.  The architecture and optimizer were tested on B,
rather than promoted merely because they had worked for SIR.  Four initial
arms combine two settings---hidden widths $(18,18)$ at learning rate
$10^{-3}$ and $(36,36)$ at $3\times10^{-4}$---with seeds 17 and 29.
Each received 256 accepted updates.  The nominated width-18, seed-29 arm was
extended to 1,024 accepted updates in 1,051 attempts.  Training used batch
size eight, shards of four, Adam parameters $(0.9,0.999,10^{-7})$, gradient
norm cap one, conditional scale cap two, and an initial affine scale equal
to half the prior standard deviation.  The map has three IAF stages with ELU
hidden activations.  Training minimizes the reverse-KL objective
\[
 \mathcal L(\psi)=\mathbb E_{z\sim N(0,I)}
 \bigl[U_E(\mathcal T_\psi(z))
        -\log|\det D\mathcal T_\psi(z)|\bigr]+\text{constant},
\]
using B's declared finite-filter endpoint.  Thirty-two common heldout latent
draws and map inverse/derivative checks nominate a frozen map; downstream HMC
provides the sampling evidence.  This finite search nominates a map;
it does not establish an optimal architecture or statistically rank maps.

\subsection{The precision failure and its localized repair}

The first B attempt exposed a numerical failure in reset recoloring, not an
error in the deterministic transition law.  For transported particles $y_i$,
the reset uses a source covariance factor $L_y$ and a target weighted factor
$L_w$ to construct
\[
 z_i=\mu_w+L_wL_y^{-1}(y_i-\bar y).
\]
Some source covariances were sufficiently ill-conditioned that their relative
spectral margins, approximately $1.32$--$1.92\times10^{-7}$ at four inspected
endpoints, fell below the FP32 guard of approximately
$2.38\times10^{-7}$ in dimension two.  The invalid factorization fallback
could then produce a very large reported covariance error.  That flagged
fallback was not an accepted scientifically valid reset.

The optional \path{conditional_fp64} policy first executes the usual FP32
reset.  If a row fails its reset validity check, it recomputes the reset
moments, factorizations, recoloring, and their tangents in FP64, using the
already computed transported cloud, and casts the repaired row back to FP32.
Healthy rows retain their original result.  It does not add process noise,
add a covariance ridge, relax the validity threshold, or rerun the whole
particle filter in FP64.  The selected B run uses this policy; A retains the
ordinary precision policy.

The dedicated check repeated the four failing endpoints and four healthy
ones over all 50 observations.  All eight became valid.  The failing rows
required six, three, one, and one local fallbacks; the healthy rows required
none and retained bit-identical likelihoods and scores.  Against full-FP64
evaluation at these eight points, the largest relative log-likelihood
difference was $6.24\times10^{-7}$ and the largest relative score-vector
error was $4.09\times10^{-5}$.  The largest covariance-restoration error was
$1.52\times10^{-4}$, below the declared one-percent tolerance.  These are
pointwise numerical checks, not an accuracy bound over the posterior.

Because this repair changes the finite numerical target at failing rows,
the earlier B chain was not appended to the repaired run.  Training resumed
from the nominated map parameters with fresh Adam state and fresh training
draws, accepting 256 of 256 further updates in about 856 seconds.  A new map
was frozen and fresh HMC warmup and sampling followed.  A separate checkpoint
serialization repair records nonfinite diagnostic counts and a null summary
when necessary, rather than crashing while writing JSON.  It does not turn
a nonfinite endpoint into a valid one.

\subsection{Reference checks, tuning, and stopping}

The model ladder covers $T=2,10,50$ for A and $T=2,3,10,50$ for B, with
separate calibration, validation, and final datasets.  The short cases test
structural propagation and likelihood/score calculations; the final posterior
runs both use $T=50$.  A is compared with Kalman likelihoods; the reference
score diagnostic uses centered finite differences of that exact likelihood.
B uses path quadrature of orders 24 and 32 for $T\leq3$.  At $T=10,50$ it
uses bootstrap-filter references with Fisher path scores, eight replications
at each of 16,384 and 65,536 particles, reporting replication MCSE and the
change under refinement.  Those finite-particle references retain unresolved
bias.  EKF and cubature calculations provide cheap comparison points; no
paired runtime study establishes a computational advantage over them.

For transparency, Table~\ref{tab:bf-structural-ab-likelihoods} records the
historical final-scope likelihood checks before the reset precision repair.
The centered error compares changes from the generating parameter at four
probe points (that parameter, two local displacements, and the prior center):
\[
 \max_{q\in\mathcal P}\left|
 [\widehat\ell_F(q)-\widehat\ell_F(q_\star)]
 -[\ell_{\rm ref}(q)-\ell_{\rm ref}(q_\star)]\right|.
\]
Centering removes a constant log-likelihood offset; it does not remove shape
error or certify posterior agreement outside these four points.
\begin{table}[htbp]
\centering\small
\caption{Historical original-precision likelihood probes.  Values are at
the generating parameter; the last column uses all four probe points.}
\label{tab:bf-structural-ab-likelihoods}
\begin{tabular}{lrrrr}
\toprule
Model & $T$ & LEDH log likelihood & Reference & Max. centered error\\
\midrule
A & 2  & $-2.57224$ & $-2.59398$ & 0.01323\\
A & 10 & $-1.72441$ & $-1.73745$ & 0.00258\\
A & 50 & $-2.04097$ & $-1.97456$ & 0.02935\\
B & 2  & $-2.53654$ & $-2.49504$ & 0.05296\\
B & 3  & $-5.99398$ & $-5.84991$ & 0.15853\\
B & 10 & $-7.98445$ & $-8.05226$ & 0.03907\\
B & 50 & $-48.00130$ & $-48.31778$ & 0.22429\\
\bottomrule
\end{tabular}
\end{table}

The saved final-scope likelihood probes precede the conditional-precision
repair and remain labeled as such.  For example, at B's $T=50$ generating
parameter they give the $q$-score $(-1.2137,2.9454,1.1045)$, compared with
$(-1.7131,4.0065,1.2723)$ from the larger bootstrap reference.  The reference
MCSEs are $(0.0274,0.1003,0.1430)$, with additional finite-particle bias not
accounted for by those MCSEs.  The first two discrepancies are substantial
relative to the reported replication error.  This historical comparison
does not measure the repaired posterior's bias, but it prevents interpreting
successful HMC as proof of exact nonlinear likelihood or posterior accuracy.

All final particle filters use $N=4032$ and FP32.  The selected controls are
16 LEDH flow steps for A and eight for B, with transport epsilon one, five
reset iterations, and one correction iteration in both cases.  The tested
control ladder was $(16,5),(8,5),(16,10)$ for flow steps and reset iterations.
These are campaign selections, not universal defaults.  HMC coordinates and
bookkeeping use FP64, as does A's exact endpoint; B additionally has the
conditional reset repair just described.  TF32 is disabled.  The repeated
numerical kernels use stable TensorFlow graphs with XLA where enabled by the
runner; Python still performs orchestration, checkpointing, and reporting.
The recorded environment is TensorFlow 2.19.1 on an NVIDIA RTX 5080.

Four chains start from dispersed prior-coordinate locations.  Kernel tuning
targets acceptance 0.70 and permits at most four leapfrog steps.  Warmup is
adaptive: at least 2,000 and at most 10,000 transitions per chain, with a
recent 1,000-draw readiness window and $\widehat R\leq1.05$.
Retained sampling starts only after warmup, checks cumulative diagnostics
every 1,000 draws per chain, and continues to $\widehat R\leq1.01$ for every
parameter or to the cap of 20,000 retained draws per chain.  Diagnostics
include physical as well as transformed coordinates.  ESS and finite
Hamiltonian energy errors are reported as explanatory diagnostics, not extra
sampling stop conditions.  The legacy storage name \emph{verification}
denotes these retained inference samples; it is not a separate validation phase.

\subsection{Posterior results and their interpretation}

Table~\ref{tab:bf-structural-ab-sampling} gives the terminal diagnostics.
Every monitored parameter has $\widehat R<1.01$.  B's fourth chain has
acceptance $76.8\%$, slightly above the preferred $65$--$75\%$ band; its pooled
acceptance is $73.075\%$.  It would therefore be wrong to report that every
chain met that preferred acceptance band.  No force fallback was recorded
in either retained run.  The reported ESS and energy summaries describe
these runs and are not a statistical comparison between methods.

\begin{table}[htbp]
\centering\small
\caption{Completed $T=50$ structural-model sampling.  Draw counts are per
chain; all four chains contribute to the posterior summaries.}
\label{tab:bf-structural-ab-sampling}
\begin{tabular}{lrr}
\toprule
Quantity & Model A & Model B\\
\midrule
Warmup draws & 2,000 & 2,000\\
Retained draws & 3,000 & 1,000\\
Maximum monitored $\widehat R$ & 1.003236 & 1.003858\\
Acceptance, chain 1 (\%) & 71.30 & 70.50\\
Acceptance, chain 2 (\%) & 72.47 & 72.30\\
Acceptance, chain 3 (\%) & 70.40 & 72.70\\
Acceptance, chain 4 (\%) & 72.07 & 76.80\\
Pooled acceptance (\%) & 71.558 & 73.075\\
Minimum bulk ESS & 771 & 985\\
Minimum tail ESS & 506 & 563\\
Maximum absolute energy error & 62.46 & 13.11\\
HMC worker wall time (hours) & 2.875 & 3.552\\
\bottomrule
\end{tabular}
\end{table}

\begin{table}[htbp]
\centering\small
\caption{Physical-parameter posterior summaries.  Intervals are marginal
95\% credible intervals from retained draws.}
\label{tab:bf-structural-ab-posterior}
\begin{tabular}{llrrr}
\toprule
Model & Parameter & Generating value & Posterior mean & 95\% interval\\
\midrule
A & $\sigma$ & 0.25 & 0.1975 & $[0.0990,0.2740]$\\
A & $\tau$ & 0.15 & 0.1444 & $[0.0547,0.2448]$\\
B & $\rho$ & 0.70 & 0.6541 & $[0.3271,0.9213]$\\
B & $\sigma$ & 0.25 & 0.3192 & $[0.1686,0.4966]$\\
B & $\beta$ & 0.80 & 0.7307 & $[0.2146,1.7434]$\\
\bottomrule
\end{tabular}
\end{table}

All five generating values lie in their marginal intervals in
Table~\ref{tab:bf-structural-ab-posterior}.  Each model has only one final
inference dataset, so this observation does not establish repeated-sampling
coverage.  B's interval for $\beta$ is wide: successful mixing need not imply
precise identification.

For A, an additional deterministic calculation integrates its exact
Kalman-posterior density in $q$ on grids with 129 and 257 points per axis.
The two posterior mean errors in $q$ of the HMC sample against the finer grid are
$0.01144$ and $0.01306$.  Their declared tolerances are $0.02411$ and
$0.03872$, respectively, each calculated as the larger of one tenth of the
reference posterior standard deviation and twice the sample mean MCSE plus
the grid-refinement difference.  The grid mean changes are below
$2.6\times10^{-12}$ and boundary mass is approximately $1.8\times10^{-12}$.
Thus the stated mean-agreement check passes.  It does not test every feature
of the posterior distribution.  B has no corresponding exact-posterior
comparison: its reported draws and intervals belong to the fixed repaired
finite-LEDH target in \eqref{eq:bf-structural-campaign-endpoints}.

The completed campaign therefore supports a specific implementation result:
structural propagation in innovation coordinates, the shared score recurrence, a
localized reset precision repair, and frozen-map HMC can execute together
and satisfy the declared sampling diagnostics on these two examples.
It supplies exact-endpoint mean agreement for A and finite-target sampling
evidence for B.  Resolving B's nonlinear posterior approximation error would
require an independent posterior reference or a controlled refinement study;
neither a smaller $\widehat R$ nor higher acceptance answers that question.

\subsection{Reproducing the recorded experiment}

The following repository files define the experiment and preserve its results.
The campaign master plan is supplemented by the dated continuation and
endpoint-recovery notes.  The analysis file distinguishes original-precision
scope probes from repaired sampling, and the terminal review records the
remaining reference limits.
\begin{quote}\small\raggedright
Executable master:\par
\path{scripts/run_structural_ledh_ab.py}\par
Model and endpoint adapters:\par
\path{bayesfilter/testing/structural_ledh_models_tf.py}\par
\path{bayesfilter/testing/structural_ledh_target.py}\par
Campaign master plan:\par
\path{docs/plans/structural-ledh-model-ab-master-2026-10-23.md}\par
Result root:\par
\path{docs/plans/artifacts/structural-ledh-ab-20261023-01/}\par
Terminal results and exact-mean comparison, relative to that root:\par
\path{attempts/model_a-hmc-01/result.json}\par
\path{attempts/model_b-reset-hmc-01/result.json}\par
\path{attempts/reset-posterior-assessment-01/result.json}
\end{quote}
The archive retains raw posterior draws and checkpoints.  Its
\path{ARCHIVE.md} gives the reconstruction command for the 21 large random
design files omitted from Git; all 21 reconstructed files were checked against
their original SHA-256 hashes during documentation closeout.

Final $T=50$ dataset seeds are 26104350 for A and 26104450 for B; their fixed
particle-design seeds are the corresponding dataset seeds plus 10,000.
Manifests preserve commands, configurations, source identities, environment,
and output hashes.  The campaign ended at 10:14 UTC on 24 October 2026,
using 11.522 charged GPU-worker hours of its 48-hour allowance.  This charge
includes compilation and worker overhead and is not measured GPU utilization.
The documentation closeout reran 75 focused CPU tests with GPU devices
deliberately hidden, covering structural completion, references, the master,
reset precision, compact transport algebra, mapped diagnostics, energy
reporting, and retained-sample semantics.  Passing these tests is evidence
about the implementation; the posterior conclusions remain those supported
by the saved experiment above.
